\section{Conclusions}\label{sec:conclusions}

We surveyed the route matchers of sixteen routers in C, C++, Rust and Go, and reworked the author's
RegexMatcher library into a route matcher, RegexMatcher~v2. Its lookup walks a trie of path segments
stored in flat arrays, answers literal-only methods from an exact-match table, and never allocates.
Its front end checks patterns, handler signatures and duplicate routes when the program compiles.

Under hypotheses pre-specified before the run, and in the setting of Section~\ref{sec:host}, v2 was
faster than the fastest competitor in \HOneSuperior{} of \HOneCells{} cells. It was not slower in
one more. At the \shape{wild} shape with \SizeTen{} routes, neither claim holds against gin. v2 ran more
instructions than gin in \HTwoLost{} cells of \SizeTen{} routes. It built its table more slowly than
the median competitor in \HFourLostBuild{} cells with large tables. The compile-time table passed its
per-cell test in \HSevenHeld{} of \HSevenCells{} cells; the rule needs all, so H7 does not hold. No
lookup of v2 allocated, every probe of the path rules passed, and the compile-time checks rejected
every malformed case. In a minimal server, the lookup difference against v1 reached the client in
every cell.

Three questions remain open. The first is whether the results hold on other processors and with other
compilers, above all for small tables. The second is why the compile-time table was slower at
\HSevenLostCell. The pre-specified run with a heap copy did not tie the slow pairs to the table's
place in memory, and it cannot rule out the placement of code. The third is what remains of the build
time for large literal tables.
